\subsection{向量测度的定义}

Ch. III, §1, No. 3 中给出的测度定义可以如下推广：

定义 1. --- 设 F 是 R 上的 Hausdorff 局部凸空间。把空间 \(\mathcal{H}(T)\) 到 F 中的每个连续线性映射都称为 T 上取值于 F 的向量测度。

定义 1 也可以表述如下：T 上取值于 F 的向量测度是一个线性映射 m （由 \(\mathcal{K}(T)\) 到 F 中），使得对 T 的每个紧子集 K ， m 在 \(\mathcal{K}(T,K)\) 上的限制对一致收敛拓扑连续。若 \(f \in \mathcal{K}(T)\) ，也用 \(\int f dm\) 或 \(\int f(t) dm(t)\) 代替 \(\mathbf{m}(f)\) 来记。取值于 R 中的测度有时称为实测度（Ch. III, §1, No. 5）或 T 上的标量测度。

例子. --- 1) \(\mathcal{H}(\mathrm{T})\) 的恒等映射是 \(\mathrm{T}\) 上取值于 \(\mathcal{H}(\mathrm{T})\) 的向量测度。

\begin{enumerate}
\def\labelenumi{\arabic{enumi})}
\setcounter{enumi}{1}
\tightlist
\item
  *设 H 是复 Hilbert 空间， \(\mathcal{L}(\mathrm{H})\) 是 H 的连续自同态所成的赋范代数。设 A 是 \(\mathcal{L}(\mathrm{H})\) 的一个子代数，它交换、闭、自伴且包含单位元；可以证明，存在一个紧空间 X 以及赋范代数 A 到代数 \(\mathcal{K}_{\mathbf{C}}(\mathrm{X}) = \mathcal{C}_{\mathbf{C}}(\mathrm{X})\) （ X 上连续复函数所成的代数）上的一个同构，该代数赋以范数 \(\| f \| = \sup_{x \in \mathrm{X}} |f(x)|\) 。逆同构在 \(\mathcal{K}(\mathrm{X})\) 上的限制是一个向量测度 \(\mathbf{m}\) （ X 上的，取值于 \(\mathcal{L}(\mathrm{H})\) ），使得 \(\mathbf{m}(fg) = \mathbf{m}(f)\mathbf{m}(g)\) 。*
\end{enumerate}

注记. --- 1) 为使线性映射 m （由 \(\mathcal{K}(\mathrm{T})\) 到 F 中）是向量测度，必须且只需：对 T 的每个紧子集 K ， m 下单位球 \(\|f\|\leqslant1\) （ \(\mathcal{K}(\mathrm{T},\mathrm{K})\) 的）的像在 F 中有界。因此，取值于 F 的向量测度这一概念，对 F 上具有相同有界集的所有 Hausdorff 局部凸拓扑都是一样的，特别地，对与 F 和 F' 之间的对偶相容的所有拓扑都是一样的（TVS, IV, §1, No. 1, 命题 1）。

\begin{enumerate}
\def\labelenumi{\arabic{enumi})}
\setcounter{enumi}{1}
\item
  设 \(T_{1}\) 为局部紧空间， \(F_{1}\) 为 R 上的 Hausdorff 局部凸空间， u 为 \(\mathcal{K}(T_{1})\) 到 \(\mathcal{K}(T)\) 中的连续线性映射， v 为 F 到 \(F_{1}\) 中的连续线性映射。若 m 是 T 上取值于 F 的向量测度，则 \(v \circ m \circ u\) 是 \(T_{1}\) 上取值于 \(F_{1}\) 的向量测度。特别地，若 g 是 T 上的连续有限数值函数，则 \(f \mapsto \mathbf{m}(gf)\) 是 T 上取值于 F 的向量测度，记作 \(g \cdot m\) ；若 h 是 T 上另一个连续有限数值函数，则 \(g \cdot (h \cdot \mathbf{m}) = (gh) \cdot \mathbf{m}\) 。
\item
  由于空间 \(\mathcal{H}(\mathrm{T})\) 是诸 Banach 空间 \(\mathcal{H}(\mathrm{T},\mathrm{K})\) 的直接极限，因而特别是桶式的（TVS, III, §4, No. 1, 命题 2 的推论及命题 3 的推论 3），为使线性映射 \(\mathbf{m}\) （由 \(\mathcal{H}(\mathrm{T})\) 到 \(\mathbf{F}\) 中）是向量测度，必须且只需：对每个 \(\mathbf{z}' \in \mathbf{F}'\) ， \(\mathbf{z}' \circ \mathbf{m}\) 是 \(\mathrm{T}\) 上的标量测度（TVS, II, §6, No. 4, 命题 5 以及 IV, §1, No. 3, 命题 7）。
\item
  鉴于注记 1，Ch. III, §2, No. 1 的命题 1 及其证明对向量测度仍然有效。因此可以把向量测度 \(\mathbf{m}\) （ \(T\) 上的）的支集定义为最大的开集 \(U \subset T\) （使 \(\mathbf{m}\) 在 \(U\) 上的限制为零者）的余集。
\end{enumerate}

